\section{Original positive companion boundary thresholds are globally analytic}

The actual companion starts at a point's literal nonmonomial-order history
year. Its ``original'' boundary therefore includes precisely those genuine
original and exceptional members already present at that year. The historical
boundary list is not fixed over the whole final manifold. Its actual
point-dependent indexed count nevertheless has global analytic thresholds,
because it is an internally selected prefix of the same actual chronological
final boundary and the real history-year bounds are analytic.

This identifies a family of genuine positive thresholds of the SAME constructed
canonical invariant, rather than assigning a different invariant to the original
input. These are the full companion-boundary-count thresholds with the true
minimum lower value. Arbitrary later lower thresholds, special/contact
presentations and the auxiliary Hilbert--Samuel chain remain separate producers.
The historical convention agrees with Bierstone--Milman\cite{bierstone-milman08-functoriality},
Section 7, Step II.B, and the recursive historical boundary counts in
Bierstone--Milman\cite{bierstone-milman97-invent-math}, Section 6.

\begin{lemma}[The true minimum lower value makes the whole Step One threshold exactly its actual level]
\label{hir:original-companion-minimum-lower-threshold-level-dictionary}
\lean{Hironaka.HironakaOriginalCompanionChronologicalBoundaryAnalytic.minimumValTy,Hironaka.HironakaOriginalCompanionChronologicalBoundaryAnalytic.minimumValTy_le,Hironaka.HironakaOriginalCompanionChronologicalBoundaryAnalytic.stepOneMinimumThreshold_iff_level}
\leanok
Let $b_k$ be the actual minimum of the literal lower value type: $0$ at
induction dimension zero and $(0,\operatorname{inl}\bot)$ at a successor.
Every actual masked Step One value satisfies
\[
 (s,b_k)\le \overline{\operatorname{inv}}(y)
 \quad\Longleftrightarrow\quad s\le\operatorname{level}(y).
\]
The statement includes special points and every later lower/contact branch.
\end{lemma}
\begin{proof}
At dimension zero, $0$ is the minimum natural number. At a successor, either
the first natural entry is positive, giving strict lexicographic inequality,
or it is zero, where the left tagged bottom is at most every actual second
entry. Thus $b_k$ is the actual minimum, and its finite lift is at most every
finite or top lower value. The literal Step One first entry is exactly its
actual original-boundary level. Lexicographic comparison with $(s,b_k)$
is therefore true precisely when that level is at least $s$, since the
second comparison is automatic at equality. The special top value and
all nonspecial finite values obey the same minimum comparison.
\end{proof}

\begin{lemma}[The actual frozen companion boundary is the genuine chronological final prefix]
\label{hir:original-companion-frozen-chronological-boundary-dictionary}
\lean{Hironaka.HironakaOriginalCompanionChronologicalBoundaryAnalytic.chronologicalMembers,Hironaka.HironakaOriginalCompanionChronologicalBoundaryAnalytic.actualCompanionLevel_eq_chronologicalBoundaryCount}
\leanok
\uses{hir:tower-split,bcalc:nonmonomial-order-tail,hir:step-one-value,bcalc:member-intersection-tower,hir:open-restriction-marked}
For an original admissible tower $T$ of length $t$ with original boundary
length $e$, and any actual split with $r$ final stages, put
\[
 E^{[r]}_T=\operatorname{take}_{e+(t-r)}E_T.
\]
The original-boundary list of the actual companion formed on its true order
mask at this split, along its actual restricted tail, is precisely the
pullback of $E^{[r]}_T$. Consequently its actual Step One level at a point of
the real tail mask equals the indexed boundary count of $E^{[r]}_T$ at that
point's same underlying original final point.
\end{lemma}
\begin{proof}
Every genuine blowup appends its actual exceptional member after the strict
transforms of the preceding members; empty members are retained at their true
chronological positions. Therefore the prefix boundary length is exactly
$e+(t-r)$, including literal empty steps. The original members at the final
stage of the tail are the first this many members of its actual final
boundary. The entire original transform dictionary identifies that boundary
with $E_T$.

The actual companion construction preserves the entire prefix boundary, as
does its actual homogenized/coefficient marked ideal used by Step One.
Transforming the boundary only depends on that same original boundary list.
The original whole open-restriction transform dictionary therefore identifies
its actual final boundary with the pullback of $E_T$ on the true tail mask.
Taking the actual original-length prefix commutes with that list pullback.
This proves the entire list identity. The indexed boundary count is unchanged
by pullback: an index passes through the restricted point precisely when its
same original member passes through the underlying point. This gives the
stated actual level identity without a caller boundary or value dictionary.
\end{proof}

\begin{lemma}[The whole actual companion history count has a finite threshold dictionary]
\label{hir:original-companion-chronological-count-finite-union-dictionary}
\lean{Hironaka.HironakaOriginalCompanionChronologicalBoundaryAnalytic.companionChronologicalCount,Hironaka.HironakaOriginalCompanionChronologicalBoundaryAnalytic.companionChronologicalCount_superlevel_eq_biUnion}
\leanok
\uses{hir:original-companion-history-year-threshold,hir:original-companion-frozen-chronological-boundary-dictionary}
For the SAME original actual history, put
$c_q(y)=\operatorname{count}(E^{[R_q(y)]}_T,y)$ and
\[
 C_{q,s}=\{y\in\operatorname{cosupp}M_T:q\le\nu_T(y),\ s\le c_q(y)\}.
\]
Then the exact entire identity is
\[
 C_{q,s}=\bigcup_{0\le r\le t}
 \bigl(L_{q,r}(T)\cap\{y:s\le\operatorname{count}(E^{[r]}_T,y)\}\bigr).
\]
No history year or selected prefix is supplied as additional data.
\end{lemma}
\begin{proof}
The literal run is at most the genuine finite tower length. A point of the
left set therefore supplies its actual $r=R_q(y)$ in the true finite union;
the selected prefix and count are exactly its own. Conversely membership
in the $r$ summand says $R_q(y)\le r$. Hence $t-r\le t-R_q(y)$ and
$E^{[r]}_T$ is an actual indexed prefix of $E^{[R_q(y)]}_T$.
The injection of the smaller prefix's indices into the larger one preserves
each member and membership through the same point. The smaller indexed
count is at most the larger indexed count, so its lower threshold gives
the actual point-dependent count threshold. Both directions retain the
same genuine current cosupport and current-order lower bound. This proves
exactly the entire displayed identity, including zero or excessive count
thresholds and all real empty chronological members.
\end{proof}

\begin{theorem}[Actual chronological companion-count thresholds are genuinely global analytic sets]
\label{hir:original-companion-chronological-count-global-analytic-producer}
\lean{Hironaka.HironakaOriginalCompanionChronologicalBoundaryAnalytic.isComplexAnalyticSubset_companionChronologicalCount_superlevel}
\leanok
\uses{hir:original-companion-chronological-count-finite-union-dictionary,hir:original-companion-history-year-global-analytic-producer,hir:boundary-count-global-analytic-superlevels}
In positive ambient dimension, for only an original marked-domain input and its literal constructed
canonical finite history, each $C_{q,s}$ is globally complex analytic,
in arbitrary induction dimension without compactness or a curve envelope.
\end{theorem}
\begin{proof}
Original canonical-domain transport internally supplies the actual final
SNC boundary $E_T$. Every true indexed prefix $E^{[r]}_T$ inherits that
same SNC structure, so its finite intersection dictionary supplies actual
holomorphic equations for each indexed-count threshold. The actual
history-year producer supplies genuine global equations for $L_{q,r}(T)$.
Intersect those actual finite presentations in every true summand of the
preceding dictionary and take their real finite union. The resulting
holomorphic equations present exactly $C_{q,s}$ on the entire original
final manifold. No global finite component family or analytic target
certificate is assumed.
\end{proof}

\begin{theorem}[Literal positive canonical companion minimum-tail thresholds have the exact whole original dictionary]
\label{hir:original-positive-canonical-companion-count-threshold-whole-dictionary}
\lean{Hironaka.HironakaOriginalCompanionChronologicalBoundaryAnalytic.positiveCanonicalCompanionCountThreshold,Hironaka.HironakaOriginalCompanionChronologicalBoundaryAnalytic.canonical_positive_companion_count_superlevel_eq_union}
\leanok
\uses{hir:original-companion-minimum-lower-threshold-level-dictionary,hir:original-companion-frozen-chronological-boundary-dictionary,hir:original-companion-chronological-count-finite-union-dictionary,hir:general-value,hir:general-canonical-tail}
For original marked-domain input and its SAME literal canonical finite
history, and every positive $q$ and indexed count threshold $s$, let
\[
 \alpha_{q,s}=(q,\operatorname{inr}(s,b_k))
\]
be the genuine tagged threshold in the literal canonical value type at
induction dimension $k+1$. Then
\[
 \{y\in\operatorname{cosupp}M_T:
       \alpha_{q,s}\le\operatorname{inv}_{k+1}(M,T,y)\}
 =\bigl(\operatorname{cosupp}M_T\cap\{\nu_T\ge q+1\}\bigr)\cup C_{q,s}.
\]
The left side is the original full value, with its actual original companion
and point-dependent historical split; it is not a redefined count invariant.
\end{theorem}
\begin{proof}
At a current cosupport point of order below $q$, the original first entry
makes the comparison false, and both right terms fail by their genuine
order bounds. At order above $q$, the first entry gives strict inequality;
the actual order-$q+1$ baseline includes the point regardless of its history
count. At order exactly $q$, positivity and the genuine original canonical
history internally supply the actual companion order mask, whole companion
marked domain, true maximal order and canonical restricted tail. Its lower
value is finite and equals the literal masked Step One value at the true
restricted point. Therefore tagged lexicographic comparison is exactly
comparison of that actual Step One value with $(s,b_k)$. The first lemma
turns this into the actual companion level threshold; the whole frozen
boundary dictionary identifies that level with $c_q(y)$. This is precisely
the right count term. All cases use actual original data; no companion
value, mask-membership or chronological dictionary is an input.
\end{proof}

\begin{theorem}[Only original inputs construct global analytic positive companion-count superlevels]
\label{hir:original-positive-canonical-companion-count-global-analytic-producer}
\lean{Hironaka.HironakaOriginalCompanionChronologicalBoundaryAnalytic.isComplexAnalyticSubset_canonical_positive_companion_count_superlevel}
\leanok
\uses{hir:original-positive-canonical-companion-count-threshold-whole-dictionary,hir:original-companion-chronological-count-global-analytic-producer,hir:canonical-first-entry-analytic-superlevel}
Every literal original positive threshold $\alpha_{q,s}$ above has a genuine
global analytic superlevel in arbitrary induction dimension. No compactness,
curve, presentation, threshold-analyticity or selected-history certificate
is supplied by the caller.
\end{theorem}
\begin{proof}
The actual coherent nonmonomial ideal supplies the whole order-$q+1$ baseline
and original coherence supplies the whole current cosupport. The chronological
companion-count producer supplies genuine analytic equations for $C_{q,s}$.
Take their finite union through the preceding exact whole original value
dictionary. Its actual holomorphic presentations define precisely the
literal canonical superlevel. This proves the positive companion boundary
layer, while every later arbitrary recursive lower threshold still requires
its own genuine semicoherent marked/contact presentation construction.
\end{proof}
